\appendix
\section{Cluster algebras}

\label{sec:groc}

Here we collect basic def\/initions for cluster algebras
to f\/ix the convention and notation,
mainly following~\cite{Fomin07}.
For further necessary def\/initions and information
for cluster algebras, see~\cite{Fomin07}.

Let $I$ be a f\/inite index set throughout the appendix.

{\bf (i) Semif\/ield.}
A {\em semifield\/} $(\mathbb{P},\oplus, \cdot)$
is an Abelian multiplicative group
endowed with a binary
operation of addition $\oplus$ which is commutative,
associative, and distributive with respect to the
multiplication in $\mathbb{P}$.
The following three examples are important.


(a) {\em  Trivial semifield.}
The {\em trivial semifield\/}
$\mathbf{1}=\{\mathrm{1}\}$
consists of the multiplicative identity element $1$
with $1\oplus 1 = 1$.

\par
(b) {\em Universal semifield.}
For an $I$-tuple of variables $y=(y_i)_{i\in I}$,
the {\em universal semifield\/}
$\mathbb{P}_{\mathrm{univ}}(y)$
consists of  all the rational functions
of the form
$P(y)/Q(y)$  (subtraction-free  rational expressions),
where $P(y)$ and $Q(y)$ are the nonzero polynomials
 in $y_i$'s with {\em nonnegative\/} integer coef\/f\/icients.
The multiplication and the addition
are given by the usual ones of rational functions.


(c) {\em Tropical semifield.}
For an $I$-tuple of variables $y=(y_i)_{i\in I}$,
the {\em tropical semifield\/}
$\mathbb{P}_{\mathrm{trop}}(y)$
is the Abelian multiplicative group freely generated by
the variables  $y_i$'s endowed with the addition~$\oplus$
\begin{gather}
\label{eq:trop}
\prod_i y_i^{a_i}\oplus
\prod_i y_i^{b_i}
=
\prod_i y_i^{\min(a_i,b_i)}.
\end{gather}



{\bf (ii)  Mutations of matrix and quiver.}
An integer matrix
$B=(B_{ij})_{i,j\in I}$  is {\em skew-sym\-met\-ri\-zab\-le\/}
if there is a diagonal matrix $D=\mathrm{diag}
(d_i)_{i\in I}$ with $d_i\in \mathbb{N}$
such that $DB$ is skew-symmetric.
For a skew-symmetrizable matrix $B$ and
$k\in I$, another matrix $B'=\mu_k(B)$,
called the {\em mutation of $B$ at $k$\/}, is def\/ined by
\begin{gather}
\label{eq:Bmut}
B'_{ij}=
\begin{cases}
-B_{ij},& \mbox{$i=k$ or $j=k$},\\
B_{ij}+\frac{1}{2}
(|B_{ik}|B_{kj} + B_{ik}|B_{kj}|),
&\mbox{otherwise}.
\end{cases}
\end{gather}
The matrix $\mu_k(B)$ is also skew-symmetrizable.

It is standard to represent
a {\em skew-symmetric} (integer) matrix $B=(B_{ij})_{i,j\in I}$
by a {\em quiver $Q$
without loops or $2$-cycles}.
The set of the vertices of~$Q$ is given by~$I$,
and we put $B_{ij}$ arrows from~$i$ to~$j$
if $B_{ij}>0$.
The mutation $Q'=\mu_k(Q)$ of quiver $Q$ is given by the following
rule:
For each pair of an incoming arrow $i\rightarrow k$
and an outgoing arrow $k\rightarrow j$ in $Q$,
add a new arrow $i\rightarrow j$.
Then, remove a maximal set of pairwise disjoint 2-cycles.
Finally, reverse all arrows incident with $k$.

{\bf (iii)  Exchange relation of coef\/f\/icient tuple.}
Let $\mathbb{P}$ be a given semif\/ield.
For an $I$-tuple $y=(y_i)_{i\in I}$, $y_i\in \mathbb{P}$
and $k\in I$, another $I$-tuple $y'$ is def\/ined
by the {\em exchange relation}
\begin{gather}
\label{eq:coef}
y'_i =
\begin{cases}
\displaystyle
{y_k}{}^{-1}, &i=k,\\
\displaystyle
y_i \left(\frac{y_k}{1\oplus {y_k}}\right)^{B_{ki}},&
i\neq k,\ B_{ki}\geq 0,\\
y_i (1\oplus y_k)^{-B_{ki}},&
i\neq k,\ B_{ki}\leq 0.
\end{cases}
\end{gather}

{\bf (iv)   Exchange relation of cluster.}
Let $\mathbb{QP}$ be the  quotient f\/ield of the group ring
$\mathbb{Z}\mathbb{P}$ of $\mathbb{P}$,
 and let $\mathbb{QP}(z)$ be the rational function f\/ield of
algebraically independent variables, say, $z=(z_i)_{i\in I}$
over $\mathbb{QP}$.
For an $I$-tuple $x=(x_i)_{i\in I}$ which
is a free generating set of $\mathbb{QP}(z)$
and $k\in I$, another $I$-tuple $x'$ is def\/ined
by the {\em exchange relation}
\begin{gather}
\label{eq:clust}
x'_i =
\begin{cases}
{x_k}, &i\neq k,\\
\displaystyle
\frac{y_k
\prod\limits_{j: B_{jk}>0} x_j^{B_{jk}}
+
\prod\limits_{j: B_{jk}<0} x_j^{-B_{jk}}
}{(1\oplus y_k)x_k},
&
i= k.
\end{cases}
\end{gather}

{\bf (v)  Seed mutation.} For the above triplet $(B,x,y)$
in (ii)--(iv), which is called a {\em seed},  the mutation
$\mu_k(B,x,y)=(B',x',y')$  at $k$ is def\/ined
 by combining \eqref{eq:Bmut}, \eqref{eq:coef},
and \eqref{eq:clust}.

{\bf (vi)  Cluster algebra}. Fix a semif\/ield $\mathbb{P}$
and a seed ({\em initial seed\/}) $(B,x,y)$, where
$x=(x_i)_{i\in I}$  are algebraically independent variables
over $\mathbb{Q}\mathbb{P}$,
and
$y=(y_i)_{i\in I}$  are elements in $\mathbb{P}$.
Starting from $(B,x,y)$, iterate mutations and collect all the
seeds $(B',x',y')$.
We call  $y'$  and $y'_i$ a {\em coefficient tuple} and
a {\em coefficient}, respectively.
We call  $x'$  and $x'_i\in \mathbb{Q}\mathbb{P}(x)$, a {\em cluster} and
a {\em cluster variable}, respectively.
The {\em cluster algebra $\mathcal{A}(B,x,y)$ with
coefficients in $\mathbb{P}$} is a
$\mathbb{Z}\mathbb{P}$-subalgebra of the
rational function f\/ield $\mathbb{Q}\mathbb{P}(x)$
generated by all the cluster variables.


{\bf (vii) $\boldsymbol{F}$-polynomial.}
The cluster algebra $\mathcal{A}(B,x,y)$ with coef\/f\/icients in
the tropical semi\-f\/ield $\mathbb{P}_{\mathrm{trop}}(y)$ is called the
cluster algebra with the tropical coef\/f\/icients
(the principal coef\/f\/icients in~\cite{Fomin07}).
There, each cluster variable $x'_i$ is an element
in $\mathbb{Z}[x^{\pm 1},y]$.
The $F$-polynomial $F'_i(y)\in \mathbb{Z}[y]$ (for $x'_i$)
is def\/ined as the specialization of $x'_i$ with $x_i=1$ ($i\in I$).

